% Generated draft: requires mathematical and source review.
\documentclass[11pt]{article}
\usepackage[margin=25mm]{geometry}
\usepackage{fontspec}
\setmainfont{Latin Modern Roman}
\usepackage{amsmath,amssymb,mathtools,unicode-math}
\setmathfont{Latin Modern Math}
\usepackage{enumitem,needspace,xurl,hyperref}
\setlength{\parindent}{0pt}
\setlength{\parskip}{5pt}
\setlength{\emergencystretch}{3em}
\newcommand{\N}{\mathbb{N}}
\newcommand{\Z}{\mathbb{Z}}
\newcommand{\Q}{\mathbb{Q}}
\newcommand{\R}{\mathbb{R}}
\newcommand{\C}{\mathbb{C}}
\newcommand{\F}{\mathbb{F}}
\newcommand{\A}{\mathbb{A}}
\newcommand{\PP}{\mathbb{P}}
\newcommand{\E}{\mathbb{E}}
\newcommand{\eps}{\varepsilon}
\newcommand{\dd}{\,\mathrm{d}}
\begin{document}
\section*{UnsolvedMath ID: 30002350}
\textbf{Source problem code:} catalogue entry
\subsection{Definitions and mathematical statement}
Let \(p\) be an odd prime and let \(n\in\Z_{\geq 1}\). Fix the notation
\[
G=D_{p^n}=\langle r,s\mid r^{p^n}=s^2=1,\ srs^{-1}=r^{-1}\rangle,
\qquad |G|=2p^n.
\]
The supplied question asks whether every local Galois extension with group \(G\) lifts to characteristic zero, calling this the Oort-group property.

The record does not specify the residue field, the precise category of local extensions, or the permitted characteristic-zero coefficient rings. Under the conventional local-action interpretation described below, the intended statement can be made precise; these conventions require confirmation against the source.

Let \(k\) be an algebraically closed field of characteristic \(p\), and let \(k[[z]]\) be the formal power series ring with its \(z\)-adic topology. A local \(G\)-action is an injective homomorphism
\[
\rho:G\longrightarrow\operatorname{Aut}^{\mathrm{cont}}_k(k[[z]]).
\]
It represents the local Galois extension with invariant subring \(k[[z]]^G\), or the corresponding extension of fraction fields. A lift of \(\rho\) consists of a complete discrete valuation ring \(R\) of characteristic zero, an identification \(R/\mathfrak m_R\cong k\), and an injective homomorphism
\[
\widetilde\rho:G\longrightarrow\operatorname{Aut}^{\mathrm{cont}}_R(R[[Z]])
\]
whose reduction modulo \(\mathfrak m_R\), under \(R[[Z]]/\mathfrak m_RR[[Z]]\cong k[[z]]\), equals \(\rho\). Here continuity on \(R[[Z]]\) refers to the \((\mathfrak m_R,Z)\)-adic topology, and \(R\) may depend on the action.

With these explicitly stated conventions, the question is: for every odd prime \(p\), every integer \(n\geq1\), every algebraically closed field \(k\) of characteristic \(p\), and every local \(D_{p^n}\)-action \(\rho\) on \(k[[z]]\), does such a lift \((R,\widetilde\rho)\) exist? This formulation addresses the local lifting property expressly described in the record. The scope of the term ``Oort group,'' including any intended global lifting formulation, remains to be checked against the cited report.
\subsection{Short English statement}
For every odd prime \(p\) and integer \(n\geq1\), must every local Galois extension with dihedral group \(D_{p^n}\) of order \(2p^n\) lift to characteristic zero? The displayed formulation uses explicitly stated local-action conventions that require source confirmation. The record labels the problem open; that label is supplied curation, not an independently verified conclusion.
\subsection{Sources}
\begin{itemize}
\item Supplied local UnsolvedMath registry record, permanent ID \(30002350\), title ``Dihedral Groups as Oort Groups,'' source problem code \texttt{catalogue entry}, in the Oberwolfach Reports Open Problems collection. The record gives proposed year 2013 and identifies no individual proposer.
\item Supplied report citation: \emph{The Arithmetic of Fields} (2014), Oberwolfach Reports, DOI 10.4171/owr/2013/30, pp. 1771--1818. \url{https://doi.org/10.4171/owr/2013/30}. The background describes Oberwolfach Reports as documenting mathematical workshops and their open-problem discussions. No problem-specific page or individual attribution is supplied.
\item Obus, \emph{A generalization of the Oort conjecture}. \url{https://arxiv.org/abs/1502.07623}.
\item \emph{The local lifting problem for D4}. \url{https://arxiv.org/abs/1706.03751}.
\item Pop, \emph{The Oort conjecture on lifting covers of curves}. \url{https://annals.math.princeton.edu/2014/180-1/p06}.
\item Supplied literature assessment, checked 2026-08-21: the record labels the problem open, reports that cyclic Oort groups and several noncyclic cases, including \(D_4\), are known, and states that the classification of all Oort groups and the conjectured odd-prime dihedral family remain unsettled. This is supplied dated curation derived from the OWR material, not independently verified here; its time-sensitive status and cited claims require independent verification.
\end{itemize}
\end{document}
